\section{Application Interface}

The coding theorem is stated over the complete received-word space and does not
depend on a particular source of observations.  An application may nevertheless
inherit the result through a distance-aware map.

Let \(O\) be a type of raw presented observations, let
\(G\subseteq O\) be the exact/valid locus, and let
\[
  \pi:O\longrightarrow(\F^2)^I
\]
be the extracted pair word.  A source distance
\(\rho_O:O\to[0,1]\) is meaningful only when \(O\) includes arbitrary or noisy
presentations; if every element of \(O\) already contains a validity witness,
then the proximity problem may collapse to the exact code locus.

The Lean companion packages this interface as a source validity model and a
source-to-code calibration.  Its only quantitative field is the pointwise
inequality
\[
  \rho_O(o)\le D_2(\pi(o)).
\]
It yields, uniformly in \(\delta\),
\[
  \rho_O(o)\ge\delta
  \Longrightarrow
  D_2(\pi(o))\ge\delta.
\]

Under this comparison, the proved universal list theorem applies to
\(\pi(o)\); the proved pair-MCA theorem applies through the same interface
once the application fixes its pair before sampling the challenge.

There are two standard ways to discharge the comparison.  Defining
\(\rho_O=D_2\circ\pi\) makes it reflexive.  A genuinely pre-existing source
metric instead requires a mathematical proof of
\(\rho_O\le D_2\circ\pi\).  Exact witnesses may certify membership in \(G\),
but the carrier \(O\) must still include arbitrary received observations.

Protocol-specific observation formats, commitments, and runtime verification
are intentionally outside the proof of
\cref{thm:capacity-grid-exact,thm:next-shell-lower,thm:pair-mca}.
Such developments may cite the two coding theorems through the interface
above; neither theorem assumes that every received word comes from an
application-specific witness space.
